\section*{Chapitre \ref{ch:nw:the-european-map} --- La carte européenne}

\begin{solution}{exo:nw:the-european-map:1}
De \qty{377.2}{\micro\second} (LD4 vers Basildon, \qty{77.4}{\kilo\metre}) à \qty{4838.7}{\micro\second} (LD4 vers Bergamo,
\qty{992.2}{\kilo\metre})~: \qty{4.46}{\milli\second} de plus, dans un sens.
\end{solution}

\begin{solution}{exo:nw:the-european-map:2}
Dans un sens, \qty{2.32}{\milli\second}~; le plancher dans l'air sur \qty{677.7}{\kilo\metre} vaut \qty{2.261}{\milli\second}~; $2.320/2.261 = 1.026$,
\qty{59}{\micro\second} au-dessus du plancher, et \qty{985}{\micro\second} en dessous du plancher de la fibre en ligne droite.
\end{solution}

\begin{solution}{exo:nw:the-european-map:3}
La \omterm{def:nw:colocation-products-and-how-they-are-sold:colo}{colocation} place les équipements de la société dans le bâtiment où la plateforme apparie les ordres, à côté de ses moteurs~; un \omterm{def:nw:the-european-map:pop}{point de présence} est un point d'accès
au réseau de la plateforme dans un autre bâtiment, à partir duquel ce réseau achemine le trafic de la société vers les moteurs. La première économise la
distance~; la seconde évite seulement à la société de construire sa propre liaison longue distance.
\end{solution}

\begin{solution}{exo:nw:the-european-map:4}
Les antennes posées sur un toit se trouvent à des distances différentes de la \omterm{def:nw:colocation-products-and-how-they-are-sold:mmr}{salle d'interconnexion} (meet-me room), et le câble reliant chaque antenne à la salle donnerait sinon
à la position la plus proche un avantage de quelques nanosecondes par mètre (chapitre~9). Égaliser depuis l'arrière de chaque antenne fait que tous les fournisseurs sans fil
partent du même point dans la course du bâtiment, ce qu'exige la règle des conditions égales du chapitre~9~; sans cela, les fournisseurs
se disputeraient des positions sur le toit plutôt que des tracés plus directs.
\end{solution}

\begin{solution}{exo:nw:the-european-map:5}
Par Slough~: $678 + 992 = \qty{1670}{\kilo\metre}$, \qty{5.57}{\milli\second} dans le vide. Par Zurich~: $307 + 205 = \qty{511}{\kilo\metre}$,
\qty{1.71}{\milli\second}. Trajet direct~: \qty{498}{\kilo\metre}, \qty{1.66}{\milli\second}. Zurich coûte un détour de \qty{13.6}{\kilo\metre},
\qty{45}{\micro\second}~; Slough plus que triple le trajet.
\end{solution}

\begin{solution}{exo:nw:the-european-map:6}
Il s'agit d'une borne supérieure sur un temps à sens unique pour les ordres allant de LD4 vers IT3~: le \omterm{def:nw:the-north-american-map:factor}{facteur de tracé} dans l'air vaut au plus $4/3.311 = 1.21$, et au plus
\qty{689}{\micro\second} au-dessus du plancher, et d'au moins \qty{839}{\micro\second} de mieux que ne le serait une fibre en ligne droite. Elle ne donne pas la
valeur typique, sa distribution, le sens des données de marché, le comportement par mauvais temps (la fibre de secours) ou
l'endroit précis où l'horloge démarre et s'arrête.
\end{solution}

\begin{solution}{exo:nw:the-european-map:7}
Demander \texttt{primary} pour XPAR ce jour-là lève \texttt{LookupError}~: la table donne deux sites principaux pour ce jour-là. C'est voulu~: une plateforme apparie
ses ordres à un seul endroit, donc deux lignes signifient que la table est fausse (un chevauchement de dates, une date de fin manquante), et un choix silencieux
mettrait un chiffre faux dans chaque plancher calculé à partir d'elle.
\end{solution}

\begin{solution}{exo:nw:the-european-map:8}
La course qui compte pour un teneur de marché sur Euronext se joue entre les sociétés qui y cotent, et n'importe laquelle peut s'installer en \omterm{def:nw:colocation-products-and-how-they-are-sold:colo}{colocation} à Bergamo, y compris une
société de Londres. Ce que la migration a changé, c'est la distance entre Euronext et les autres plateformes londoniennes~: une société qui cote Euronext depuis Bergamo est
à trois millisecondes de ses couvertures londoniennes, et celle qui cote depuis Londres est à trois millisecondes du carnet d'ordres d'Euronext. Elle doit choisir où chaque
moteur se place~; elle n'est pas pour autant exclue.
\end{solution}

\begin{solution}{pb:nw:the-european-map:1}
\textbf{Partie I.}
\begin{enumerate}
\item \qty{77.4}{\kilo\metre}~: \qty{258.0}{\micro\second} dans le vide, \qty{377.2}{\micro\second} dans une fibre en ligne droite.
\item \qty{603.2}{\kilo\metre}~: \qty{2012}{\micro\second} et \qty{2942}{\micro\second}.
\item Un triangle très aplati~: 678, 603 et \qty{77}{\kilo\metre}~; Euronext se trouvait presque à Londres.
\item N'importe où dans le triangle londonien, LD4 compris~: quelques centaines de microsecondes de Basildon, des Docklands et des \omterm{def:nw:the-european-map:pop}{points de
  présence} londoniens.
\end{enumerate}
\textbf{Partie II.}
\begin{enumerate}[resume]
\item \qty{992.2}{\kilo\metre}~: \qty{3309.7}{\micro\second} dans le vide, \qty{4838.7}{\micro\second} dans la fibre.
\item $+\qty{3052}{\micro\second}$ dans le vide et $+\qty{4461}{\micro\second}$ dans la fibre.
\item \qty{497.6}{\kilo\metre}, \qty{1659.9}{\micro\second} dans le vide~: \qty{352}{\micro\second} de moins que depuis Basildon
  (\qty{515}{\micro\second} de moins dans la fibre).
\item Un triangle large~: 678, 498 et \qty{992}{\kilo\metre}~; Francfort est désormais plus proche d'Euronext que ne l'est Londres.
\end{enumerate}
\textbf{Partie III.}
\begin{enumerate}[resume]
\item $1.2 \times 4\,838.7 = \qty{5806}{\micro\second}$.
\item Moins de \qty{4000}{\micro\second}~: au moins \qty{1806}{\micro\second} gagnées.
\item Pour coter dans le carnet d'Euronext, oui~: ses annulations et ses recotations ne parcourent alors que des mètres, et non un millier de kilomètres, et Eurex est plus proche de Bergamo
  (\qty{1.66}{\milli\second}) que de LD4 (\qty{2.26}{\milli\second}). Elle perd la proximité des plateformes londoniennes (environ \qty{3.1}{\milli\second}
  depuis Bergamo), et la facilité d'accès de ses équipes à Londres~; beaucoup de sociétés exploitent les deux.
\item $306.7 + 204.5 - 497.6 = \qty{13.6}{\kilo\metre}$, \qty{45}{\micro\second} de lumière dans le vide.
\end{enumerate}
\textbf{Partie IV.}
\begin{enumerate}[resume]
\item \emph{Résultat principal}~: le plancher à sens unique entre LD4 à Slough et les moteurs d'appariement d'Euronext est passé de \qty{0.26}{\milli\second} à
  \qty{3.31}{\milli\second} dans le vide (\qty{+3.05}{\milli\second}) et de 0.38 à \qty{4.84}{\milli\second} dans la fibre
  (\qty{+4.46}{\milli\second})~; le triangle Londres--Francfort--Euronext est passé de 678, 603 et \qty{77}{\kilo\metre} à 678, 498 et
  \qty{992}{\kilo\metre}.
\item Pour que les clients londoniens puissent se raccorder à Bergamo, tester et migrer avant la bascule, sans construire leurs propres liaisons vers l'Italie.
\item Deux lignes, Basildon jusqu'au 5 juin 2022 et Bergamo à partir du 6 juin, chacune avec sa source~: une requête par date donne alors la bonne réponse pour n'importe quel
  jour, et l'historique est conservé pour les backtests qui doivent utiliser les latences de leurs propres dates.
\item Le déménagement du London Stock Exchange de la City vers les Docklands, prévu pour octobre 2022 et réalisé en février 2023.
\item Qu'un relais doive se trouver sur la géodésique Francfort--Bergamo ou à proximité~; Zurich coûte \qty{45}{\micro\second}, Londres \qty{3.9}{\milli\second}.
\item Basildon, Bergamo et Upplands Väsby. Une erreur d'un kilomètre ou deux déplace un plancher de quelques microsecondes, contre un changement de trois
  mille microsecondes~: cela ne change rien à ce résultat.
\item Une fibre dont la lumière se propage presque comme dans l'air rapprocherait le plancher de la fibre de celui de l'air, soit \qty{3.3}{\milli\second}, et supprimerait l'essentiel
  de l'avantage des liaisons hertziennes micro-ondes sur ce couloir, avec la bande passante de la fibre et son indépendance vis-à-vis de la météo (chapitre~13).
\item Le déménagement a éloigné Euronext de trois millisecondes et a obligé la société à choisir, moteur par moteur, entre être proche d'Euronext et être proche de Londres.
\end{enumerate}
\end{solution}

\begin{solution}{iq:nw:the-european-map:1}
Euronext à Bergamo (Aruba IT3), Eurex et Xetra à Francfort (Equinix FR2), le London Stock Exchange dans les Docklands, ICE Futures Europe à
Basildon, SIX à Zurich (ZH4), Nasdaq Nordic près de Stockholm. Londres (Slough) à Francfort est d'environ \qty{678}{\kilo\metre}~: \qty{2.26}{\milli\second}
à sens unique dans l'air, \qty{3.3}{\milli\second} dans une fibre en ligne droite.
\iqlookfor{Des bâtiments plutôt que des villes~; le déménagement d'Euronext de 2022~; les planchers à sens unique dans les deux milieux.}
\end{solution}

\begin{solution}{iq:nw:the-european-map:2}
Un point d'accès au réseau de la plateforme dans un bâtiment où sont déjà ses clients, pour qu'ils se connectent par une \omterm{def:nw:colocation-products-and-how-they-are-sold:mmr}{interconnexion} (cross-connect) et que la plateforme
achemine le trafic vers ses moteurs. Une bourse en ouvre une à l'étranger lorsque ses clients s'y trouvent~: Euronext en a ouvert deux à Londres avant de déménager vers
l'Italie.
\iqlookfor{PoP contre \omterm{def:nw:colocation-products-and-how-they-are-sold:colo}{colocation}~; le réseau de la plateforme et sa latence~; la résilience (deux PoP)~; l'\omterm{def:nw:colocation-products-and-how-they-are-sold:mmr}{interconnexion} (cross-connect).}
\end{solution}

\begin{solution}{iq:nw:the-european-map:3}
Le conserver sous forme de données, avec une source et une date de lecture par ligne, une ligne par plateforme, rôle et période~; revérifier les avis des plateformes selon un
calendrier~; signaler les lignes plus anciennes qu'un seuil~; ne jamais écraser une ligne quand une plateforme déménage, la fermer et en ajouter une autre.
\iqlookfor{Provenance~; historique conservé pour les backtests~; contrôles automatiques de fraîcheur~; abonnement aux avis des plateformes.}
\end{solution}

\begin{solution}{iq:nw:the-european-map:4}
Lire les documents de migration de la plateforme~; calculer les nouveaux planchers à partir de chaque site de la société~; décider où doit se trouver le moteur de chaque stratégie~;
commander l'espace, l'alimentation et la connectivité dans le nouveau bâtiment ainsi que les liaisons vers celui-ci, avec les délais d'approvisionnement~; tester lors des répétitions de la plateforme~; planifier le week-end de bascule
et le repli~; mettre à jour le registre des sites et chaque hypothèse de latence dans la recherche.
\iqlookfor{Ordre des opérations et délais d'approvisionnement~; la répétition~; hypothèses de la recherche mises à jour, pas seulement celles de la production.}
\end{solution}

\begin{solution}{iq:nw:the-european-map:5}
Parce que leurs clients sont ailleurs et que l'attrait de leur propre bâtiment dépend de la vitesse à laquelle ces clients peuvent l'atteindre~; une liaison hertzienne
micro-ondes exploitée par la plateforme, avec un accès égal (SIX de Zurich vers Francfort, Euronext de Londres vers Bergamo), vend de la vitesse selon les règles d'accès équitable de la plateforme et
apporte du flux d'ordres.
\iqlookfor{Accès équitable~; revenus~; concurrence avec les réseaux de tiers~; un sens ou les deux.}
\end{solution}

\begin{solution}{iq:nw:the-european-map:6}
Les planchers indiquent~: près de la ligne Francfort--Bergamo, ou dans l'un des deux bâtiments reliés par une liaison rapide à l'autre~; Zurich n'est qu'à
\qty{45}{\micro\second} de la ligne directe. Vérifier avec le registre et les géodésiques, puis avec les tracés publiés et leurs \omterm{def:nw:the-north-american-map:factor}{facteurs de tracé},
et enfin en mesurant une liaison candidate à l'aide de captures horodatées aux deux extrémités.
\iqlookfor{Inégalité triangulaire~; quel côté de la latence compte le plus pour la stratégie~; la mesure, et pas seulement les planchers.}
\end{solution}
